\chapter{Атрибуция общего носителя радиационного охлаждения}

Предыдущий гейт сохранил точный положительный ток полной Real-пары,
равный $4\pi^2$ за четырёхтактный цикл, но не связал его с семейным
состоянием $R$. Последний минимальный тест маршрута должен различить два
утверждения: возможность формально поместить оба сектора в тензорное
произведение и наличие уже выведенного родителя, который действительно
переносит энергию и энтропию между ними.

\section{Три имеющихся носителя}

Дискретный нелинейный родитель построен на
\begin{equation}
 \mathbb C^2_{\rm dir}\otimes E,
 \qquad E=M_{20\times15}(\mathbb C),
 \qquad \dim_{\mathbb C}E=300.
\end{equation}
Однако совместимость с полным $M_{20}$--$M_{15}$-бимодулем даёт
\begin{equation}
 \operatorname{End}_{M_{20}-M_{15}}(E)=\mathbb C I_E.
 \label{eq:v6-common-carrier-full-commutant}
\end{equation}
Следовательно, полный родитель допускает только скалярное действие на
внутренней кратности и не содержит семейно-селективной монеты.

После ограничения до наблюдаемой алгебры возникает одна семейная линия
\begin{equation}
 H_{15}=Q_L^{(6)}\oplus L_L^{(2)}\oplus u_R^{(3)}
 \oplus d_R^{(3)}\oplus e_R^{(1)}.
\end{equation}
Её пятимерный центральный коммутант различает виды частиц внутри одной
семьи. Compacton использует угол $L_L\oplus e_R$ и пространственный
фактор $\mathbb C^2_{\rm dir}$. Семейный параметр порядка
$R\in M_3(\mathbb R)$ в это пространство не входит: его три измерения
нумеруют семейную кратность, а не пять центральных блоков $H_{15}$.

Наконец, аффинный проектор $P_3$ действительно имеет каноническую
коизометрию $V$ на семейный триплет. Но ранее проверенная связь $\rho V$
действует на нём скалярно и сохраняет $R=I_3/3$. Она не содержит ни
хирального compacton-угла, ни исходящих пространственных мод.

\section{Факторизованный подъём и точный запрет}

Можно объявить общий кинематический носитель
\begin{equation}
 \mathcal H_{\rm kin}
 =\mathbb C^3_{\rm fam}\otimes
  \bigl(\mathbb C^2_{\rm dir}\otimes(L_L\oplus e_R)\bigr).
 \label{eq:v6-common-carrier-kinematic-product}
\end{equation}
Это корректное пространство, но его нет в предыдущем действии как
атрибутированного взаимодействующего угла. Единственный подъём имеющейся
compacton-динамики, не вводящий нового оператора, равен
\begin{equation}
 U_{\rm can}=I_3\otimes U_{\rm walk}.
\end{equation}

\begin{theorem}[Факторизованная радиация не меняет семейное состояние]
Для произвольного, в том числе запутанного, состояния $\Omega$ на
$\mathcal H_{\rm kin}$ выполняется
\begin{equation}
 \operatorname{Tr}_{\rm walk}
 \left[(I_3\otimes U_{\rm walk})\Omega
 (I_3\otimes U_{\rm walk}^{*})\right]
 =\operatorname{Tr}_{\rm walk}\Omega.
 \label{eq:v6-common-carrier-partial-trace-no-go}
\end{equation}
Следовательно, спектр $R$, его энтропия и параметр порядка $Q=R-I_3/3$
сохраняются точно. Ненулевой ток $4\pi^2$ остаётся током сектора
блуждания и
не является тепловым током семейного состояния.
\end{theorem}

Равенство \eqref{eq:v6-common-carrier-partial-trace-no-go} следует из
цикличности частичного следа по фактору, на котором действует унитарный
оператор. Машинный аудит проверяет его на двадцати четырёх случайных
положительных совместных состояниях; максимальный остаток меньше
$10^{-14}$. Тот же аудит восстанавливает одномерный коммутант полного
$M_3(\mathbb C)$-действия на семейном триплете, как требует лемма Шура.

\section{Что потребовалось бы для ненулевого переноса}

Изменение $R$ требует оператора, не факторизующегося как тождество на
семейном факторе, например
\begin{equation}
 \sum_a T_a(Q)\otimes J_a^{\rm out}.
 \label{eq:v6-common-carrier-missing-interaction}
\end{equation}
Но ни $T_a(Q)$, ни его коэффициент, ни общая следовая нормировка, ни
разбиение на систему и резервуар в проекте не выведены. Зависимость от
заранее выбранного проектора $P$ была бы круговой: она предполагала бы
ось, которую механизм должен породить. Общая теория подсистем и
безошибочных факторов подчёркивает то же различие между кратностью,
коммутантом и взаимодействием
\cite{common-carrier-zanardi1997,common-carrier-knill2000}.

\begin{nogocriterion}[Совместное вложение не является связью]
Нельзя считать \eqref{eq:v6-common-carrier-kinematic-product} физическим
родителем лишь потому, что тензорное произведение математически
существует. Для охлаждения необходим уже атрибутированный нескалярный
переплетающий оператор, меняющий спектр $R$, и единая энергетическая
нормировка.
Добавление $Q$- или $P$-зависимой стрелки с новым коэффициентом является
новой моделью, а не закрытием текущего маршрута.
\end{nogocriterion}

\begin{proposition}[Итог атрибуции]
Полный $M_{300}$-фон, физический $H_{15}$-угол и аффинный семейный
триплет покрывают нужные компоненты только по отдельности. Первый имеет
скалярный коммутант, второй не содержит семейный $R$, третий не содержит
пространственно-хиральную радиацию и действует изотропно. Поэтому
канонического общего взаимодействующего носителя в закрытом багаже нет.
Радиационно-охлаждающий маршрут закрывается на уровне существующего
родителя, тогда как точный коэффициент $4\pi^2$ сохраняется как отдельный
динамический инвариант compacton-блуждания.
\end{proposition}

\section{Следующий гейт}

Следующий документ
\path{version6_spectral_transition_post_radiative_bridge_final_dynamic_status_gate}
должен свести финальный динамический ledger Тома VI: отдельно записать
доказанное рождение Real-пары и проекторную фазовую кинематику, отдельно
--- отсутствие выведенного самозапускающегося охлаждающего закона.

Нулевая гипотеза: после закрытия compacton- и радиационного маршрутов в
Томе VI не осталось незаписанного коэффициент-свободного механизма,
который меняет $R$. Стоп-критерий: не открывать новую динамическую ветвь
без нового родительского действия, общей нормировки и слепого наблюдаемого
следствия.

\begin{protocol}
\begin{enumerate}
 \item общий кинематический тензорный носитель: \textbf{можно объявить};
 \item он уже выведен как один взаимодействующий родитель: \textbf{нет};
 \item коммутант полного бимодуля: \textbf{$\mathbb C I$};
 \item коммутант полного семейного $M_3$-действия: \textbf{одномерный};
 \item $H_{15}$ содержит compacton-угол: \textbf{да};
 \item $H_{15}$ содержит семейный параметр $R$: \textbf{нет};
 \item аффинный $P_3$ содержит семейный триплет: \textbf{да};
 \item каноническая аффинная связь выбирает ось: \textbf{нет};
 \item остаток изменения $R$ при $I_3\otimes U_{\rm walk}$:
 \textbf{меньше $10^{-14}$};
 \item новый коэффициент-свободный переплетающий оператор:
 \textbf{не найден};
 \item ток $4\pi^2$: \textbf{сохранён, но не атрибутирован к $R$};
 \item радиационное охлаждение существующим родителем: \textbf{закрыто}.
\end{enumerate}
\end{protocol}

Машинный сертификат сохранён в
\path{s2t_v6_spectral_transition_radiative_cooling_common_carrier_attribution_gate_results.json}.

\begin{thebibliography}{9}
\bibitem{common-carrier-zanardi1997}
P.~Zanardi and M.~Rasetti,
\emph{Noiseless Quantum Codes}, Physical Review Letters 79 (1997),
3306--3309, \path{doi:10.1103/PhysRevLett.79.3306}.

\bibitem{common-carrier-knill2000}
E.~Knill, R.~Laflamme, and L.~Viola,
\emph{Theory of Quantum Error Correction for General Noise},
Physical Review Letters 84 (2000), 2525--2528,
\path{doi:10.1103/PhysRevLett.84.2525}.
\end{thebibliography}